\section{Benchmark design}\label{sec:benchmark}

\subsection{A twin with deliberate model mismatch}\label{sec:twin}

A detector built on a physical model and evaluated on data generated by the same model would be tested against itself. The twin used here therefore contains effects that the honest model of \cref{sec:family} does not describe, so that every result is obtained under model mismatch:
\begin{enumerate}[label=(\alph*)]
\item afterpulsing with probability up to 3\,\%, which the monitor knows only through a datasheet value carrying a $\pm10$\,\% error;
\item window-to-window jitter of each intensity class with relative standard deviation up to 1\,\%, independent between signal and decoy, which the honest model cannot represent with its single ratio $\rho$;
\item polarisation drift of the misalignments, an autoregressive process with per-window innovation up to $2\times10^{-4}$;
\item temperature-driven modulation of the dark-count probabilities with amplitude up to 10\,\%;
\item a random walk and a slow periodic modulation of the channel transmittance.
\end{enumerate}
Links are drawn with lengths of 5 to 60\,km at 0.18 to 0.22\,dB/km, receiver efficiencies of 0.10 to 0.50, misalignments of 0.5 to 2\,\%, dark-count probabilities of $2\times10^{-7}$ to $5\times10^{-6}$ per gate and detector efficiency ratios of 0.85 to 1.15; \cref{tab:twin} in \cref{app:params} lists all ranges. Each link is observed for 200 windows of $N=10^7$ pulses (one second at 10\,MHz), of which the first $W=60$ are attack-free commissioning windows. Counts are sampled from the multinomial laws \eqref{eq:mult}, so every cell, including vacuum double clicks, carries its own shot noise.

\subsection{Splits, attacks and faults}\label{sec:splits}

Each seed generates disjoint populations of links. The \emph{training} split (60 honest links, plus 30 links for each known attack used only by the supervised baseline) fixes the channel standardisation of \cref{eq:combine} and trains the learned baselines. The \emph{validation} split (50 honest links and 25 links with benign loss faults) fixes the thresholds of every method. The \emph{test} split contains 50 honest links, 25 links with loss faults (steps or ramps of 0.5 to 3\,dB), 15 links with a dark-count fault (a rise by a factor 1.5 to 5), 35 links for each of the six attack families and 20 for each of three compositions (PNS with blinding, time-shift with Trojan horse, intercept--resend with RNG bias). A \emph{long-link} split of 60 to 120\,km with receiver efficiencies of 0.05 to 0.20, a regime absent from training, contains 25 honest links and 15 links for each known attack.

Attacks start uniformly between 20 and 100 windows after commissioning, and windows are labelled positive only when an attack is active, so pre-onset windows of attacked links count as honest. Attack strengths are drawn log-uniformly over physically motivated ranges (\cref{tab:attackranges} in \cref{app:params}): intercepted fractions of 0.2 to 20\,\%, PNS and blinding fractions of 1 to 100\,\%, time-shift fractions of 2 to 100\,\% with mismatch factors of 0.1 to 0.5, back-scatter probabilities of $10^{-8}$ to $3\times10^{-6}$ per gate, and RNG bias amplitudes of $5\times10^{-5}$ to $2\times10^{-3}$ with periods of 10 to 80 windows. For every attacked run we also compute $\Dstar$ at the run's mean physical state, which allows each delay to be compared with its own bound.

\subsection{Comparison methods}\label{sec:baselines}

All comparison methods receive the same information as \qtrace{}: the 19 scalar telemetry quantities a deployed node exports (gains and QBERs per intensity and basis, vacuum yield, decoy-to-signal gain ratio, double-click rates, unsifted detector balance, click rate, bit and basis balance, $Y_1^L$, $e_1^U$ and the key rate), standardised per link against the same commissioning windows, together with their trailing 16-window means. We compare
\begin{itemize}
\item three \emph{industry monitors}: the static 11\,\% QBER threshold, a one-sided CUSUM on the standardised signal QBER \cite{Page1954}, and a two-sided CUSUM on the decoy-state estimates $Y_1^L$, $e_1^U$ and the gain ratio, i.e.\ the decoy consistency check of deployed systems made sequential;
\item four \emph{unsupervised detectors} trained on honest data only: an isolation forest \cite{Liu2008IF}, a one-class SVM \cite{Scholkopf2001}, a multilayer-perceptron autoencoder, and a one-class SVM on a fidelity quantum kernel \cite{Havlicek2019,SchuldKilloran2019} with a six-qubit ZZ feature map on principal components, evaluated with exact statevectors that we verified against Qiskit \cite{Qiskit2024} to unit fidelity;
\item a \emph{supervised} gradient-boosted classifier (XGBoost \cite{ChenGuestrin2016}) trained on the three known attack families.
\end{itemize}
Standard implementations from scikit-learn \cite{Pedregosa2011} and XGBoost are used with fixed hyperparameters listed in the repository. Every method's threshold is the $1-\alpha$ quantile of its score on the same validation links, so all methods operate at the same per-window false-alarm rate.

\subsection{Metrics and protocol}\label{sec:metrics}

The primary metric is the fraction of attacked links on which a method raises an alarm after onset, at a per-window false-alarm rate of $\alpha=10^{-3}$ ($10^{-2}$ in \cref{app:results}). We also report the window-level AUROC (attack-active windows against all windows of honest test links), the detection delay in windows, the key generated between onset and first alarm (key exposed), and the fraction of honest, loss-fault and long-link honest runs with at least one alarm. All design decisions, including the horizons, the reference length and the shrinkage constants, were fixed on development seeds 0 to 3. The reported results use seeds \SeedList{}, which played no role in development; means are reported with 95\,\% Student-$t$ confidence intervals over seeds. In total the evaluation comprises \NRunsTotal{} links and \NWindowsTotal{} telemetry windows; one seed takes about \WallPerSeed{} minutes on a single CPU core.
